% Template for post or presentation abstracts submitted to ILAS 2022
% 1. Fill in the presenter's information (name, affiliation, email
% address) and talk information (title of presentation, abstract).
% 2. Compile to make sure there are no errors.
% 3. Save the .tex file for your abstract with a distinctive name,
% ideally "FamilynameGivenname.tex".
% 4. Upload your .tex file to https://clr.ie/129557
%       
% * Please keep your LaTeX commands simple, and avoid the use of
%    user-defined macros.
% * Include details of co-authors and funding acknowledgements after the abstract.
% * Upload only the LaTeX file (with .tex extension).
% * Questions: Contact Niall Madden (Niall.Madden@NUIGalway.ie)

\documentclass[a4paper]{amsart} 
\usepackage{amssymb} 
\usepackage{hyperref}
\pagestyle{empty}
\date{} 

%% IGNORE THE NEXT 4 LINES
\makeatletter % 
\def\labelenumi{\theenumi.}
\def\theenumi{\@arabic\c@enumi}
\makeatother
%% STOP IGNORING NOW

\begin{document}

\title{Self-adjoint operators associated with Hankel moment matrices \LaTeX\ abstract for a talk or poster at ILAS 2022}
\author{Christian Berg} %Presenting author only
\address{University of Copenhagen, Denmark} % Affiliation of presenting author
\email{berg@math.ku.dk} % Email address of presenting author only

\maketitle
Let $\mathcal M^*$ denote the set of positive measures with moments of any order and infinite support on the real line $\mathbb R$. The moment sequence  of $\mu\in\mathcal M^*$ is denoted
\begin{equation}\label{eq:mom}
m_n=\int_{-\infty}^\infty x^n\,d\mu(x),\quad n=0,1,\ldots,
\end{equation} 
and we let
\begin{equation}\label{eq:Han}
\mathcal H=\mathcal H_\mu=\left(m_{k+l}\right)_{k,l=0}^\infty=\{m_{k+l}\}
\end{equation}
denote the corresponding Hankel matrix.


 Denoting by $\mathcal F$ the set of sequences $g\in\ell^2$ with only finitely many non-zero entries, we have  a positive sesquilinear form $Q$ defined on $\mathcal F\times\mathcal F$
\begin{equation}\label{eq:quadr}
Q(g,h)=\sum_{k,l=0}^\infty m_{k+l}g_k\overline{h_l}, 
\end{equation}
called the Hankel form associated with the  sequence $(m_n)$.

Widom proved in 1966 that the form $(Q,\mathcal F)$ is bounded on $\ell^2$ if and only if $m_n=O(1/n)$. In 2016 Yafaev proved that the form is closable if
 $m_n=o(1)$ and characterized the closure of the form based on his earlier work on quasi-Carleman operators. We give a new proof of the description of the closure based entirely on moment considerations. In 2020 Berg and Szwarc pointed out that the form $(Q,\mathcal F)$ 
is also closable if $\mu$ is indeterminate of if $\mu$ is determinate with finite index of determinacy. 
We give  a description of the self-adjoint   operators $H=H_\mu$ (bounded or unbounded) in the Hilbert space $\ell^2$
associated with the closed Hankel forms in the  three cases mentioned, where the form is closable. 



% The abstract  ***no more than one page***


\emph{This is joint work with Ryszard Szwarc (Wroc{\l}aw), see \cite{besz8}.}


\begin{thebibliography}{1}
\bibitem{besz8} Christian Berg and Ryszard Szwarc.
\newblock Self-adjoint operators associated with Hankel moment matrices.
\newblock {\em  ArXiv:2105.07252}.


\end{thebibliography}


\end{document}


% Please leave the next bit alone  
%%% Local Variables: 
%%% mode: latex
%%% TeX-master: "../ILAS-2022"
%%% End:
